% !TeX program = xelatex
% Báo cáo ngắn gọn: 5 phương pháp + demo 2 cặp + hạn chế + kế hoạch.
% Biên dịch: xelatex -interaction=nonstopmode "[06-10]-Report.tex"
% Số liệu: results/*/metrics.json + threshold.json, ngưỡng chọn trên validation.
\documentclass[11pt,a4paper]{article}

\usepackage{fontspec}
\setmainfont{Times New Roman}
\usepackage[a4paper,top=2.0cm,bottom=2.0cm,left=2.2cm,right=1.9cm]{geometry}
\usepackage{booktabs}
\usepackage{enumitem}
\usepackage[vietnamese]{babel}
\usepackage{xcolor}
\definecolor{evyellow}{HTML}{B8860B}
\newcommand{\mline}[1]{#1}
\newcommand{\eline}[1]{\colorbox{evyellow!30}{#1}}
% Màu cặp tương ứng: hai phía cùng số là hai đoạn model cho là khớp nhau.
% Màu xoay vòng theo số cặp, cặp 8 trùng màu cặp 0.
\definecolor{pc0}{HTML}{F4B6C2}
\definecolor{pc1}{HTML}{B5D8F7}
\definecolor{pc2}{HTML}{B8E6B8}
\definecolor{pc3}{HTML}{F7E08B}
\definecolor{pc4}{HTML}{D9C2F0}
\definecolor{pc5}{HTML}{F5C99B}
\definecolor{pc6}{HTML}{A8E0E0}
\definecolor{pc7}{HTML}{E8C39E}
\newcommand{\pline}[2]{%
  \ifcase#1\colorbox{pc0}{#2}\or\colorbox{pc1}{#2}\or\colorbox{pc2}{#2}\or%
  \colorbox{pc3}{#2}\or\colorbox{pc4}{#2}\or\colorbox{pc5}{#2}\or%
  \colorbox{pc6}{#2}\or\colorbox{pc7}{#2}\fi}

\title{\vspace{-1.0cm}\textbf{Báo cáo kết quả huấn luyện + demo}\\[4pt]
       {\normalsize Đề tài: Phát hiện trùng lặp code}}
\date{\today}

\begin{document}
\maketitle
\vspace{-0.8cm}

\section{Tổng quan}
Phát hiện trùng lặp code trên POJ-104 với 5 phương pháp:
CodeBERT, GraphCodeBERT, ASTNN, CPG-GNN, Hybrid.
Ngưỡng chọn trên validation, giữ nguyên khi test và khi demo.

\section{Các phương pháp}
\begin{itemize}[leftmargin=*,itemsep=0pt,topsep=2pt,parsep=0pt]
\item \textbf{CodeBERT}: \texttt{codebert-base}, 126M tham số,
  ghép ba vector A, B và trị tuyệt đối của hiệu rồi qua MLP.
\item \textbf{GraphCodeBERT}: \texttt{graphcodebert-base}, 126M tham số,
  như CodeBERT, thêm luồng dữ liệu. CodeBERT và GraphCodeBERT cùng LR 2e-5, batch 16,
  đầu vào tối đa 256 token, 3 epoch.
\item \textbf{ASTNN}: 0,4M tham số, huấn luyện từ đầu,
  cây cú pháp qua RvNN và Bi-GRU, 7 epoch, LR 1e-3.
\item \textbf{CPG-GNN}: 0,2M tham số, huấn luyện từ đầu,
  đồ thị CPG 6 loại cạnh qua GAT 3 tầng, 4 head, hidden 128,
  5 epoch, tối đa 1500 node.
\item \textbf{Hybrid}: 128M tham số, hợp GraphCodeBERT và CPG-GNN
  ở mức embedding, warm-start từ GraphCodeBERT và CPG-GNN, fine-tune 4 epoch.
\end{itemize}

\section{Kết quả huấn luyện}
Giá trị là F1 trên validation sau mỗi epoch, đánh số từ 0.
In đậm là epoch tốt nhất, ngưỡng quyết định lấy tại đó.
{\small
\begin{tabular*}{\textwidth}{@{\extracolsep{\fill}}lccccccc@{}}
\toprule
\textbf{PP} & \textbf{E0} & \textbf{E1} & \textbf{E2} & \textbf{E3} &
\textbf{E4} & \textbf{E5} & \textbf{E6} \\
\midrule
CodeBERT & 0,919 & \textbf{0,919} & 0,912 & --- & --- & --- & --- \\
GraphCodeBERT & \textbf{0,916} & 0,914 & 0,911 & --- & --- & --- & --- \\
ASTNN & 0,789 & 0,814 & 0,841 & \textbf{0,853} & 0,804 & 0,823 & 0,841 \\
CPG-GNN & 0,780 & 0,796 & 0,800 & 0,817 & \textbf{0,819} & --- & --- \\
Hybrid & \textbf{0,914} & 0,893 & 0,859 & 0,879 & --- & --- & --- \\
\bottomrule
\end{tabular*}\par}
CodeBERT và GraphCodeBERT đạt đỉnh ngay đầu rồi đi ngang, train 3 epoch là đủ.
Tổng giờ train trên GPU Tesla T4: CodeBERT 7,0, GraphCodeBERT 6,9,
ASTNN 0,7, CPG-GNN 1,0, Hybrid 10,8.
ASTNN lên tới E3 rồi rớt ở E4, dấu hiệu overfit nên dừng sớm giữ E3.
CPG-GNN tăng đều tới epoch cuối, còn có thể train thêm.
Hybrid đỉnh ngay E0 rồi giảm, warm-start tốt nhưng fine-tune cần LR nhỏ hơn.

\section{Kết quả test}
{\small
\begin{tabular*}{\textwidth}{@{\extracolsep{\fill}}lccccc@{}}
\toprule
\textbf{PP} & \textbf{Prec.} & \textbf{Rec.} & \textbf{F1} & \textbf{AUC} & \textbf{Ngưỡng} \\
\midrule
CodeBERT      & 0,931 & 0,890 & \textbf{0,910} & \textbf{0,973} & 0,500 \\
GraphCodeBERT & 0,883 & 0,933 & 0,907 & 0,971 & 0,512 \\
ASTNN         & 0,760 & 0,894 & 0,821 & 0,889 & 0,557 \\
CPG-GNN       & 0,695 & 0,885 & 0,779 & 0,864 & 0,420 \\
Hybrid        & 0,905 & 0,922 & \textbf{0,913} & 0,970 & 0,370 \\
\bottomrule
\end{tabular*}\par}

% Mục 5 cũ là bảng điểm demo đã bỏ, điểm số chuyển vào 6.1 và 6.2.
% Giữ số mục Minh hoạ là 6 cho khớp cấu trúc đã chốt.
\setcounter{section}{5}
\section{Minh hoạ}
Cách đọc hình:
\begin{itemize}[leftmargin=*,itemsep=0pt,topsep=2pt,parsep=0pt]
\item Hai phía cùng màu là hai đoạn chính model đó cho là khớp nhau.
\item Mỗi model một bảng màu riêng vì mỗi họ nhìn code bằng tín hiệu riêng.
\item CodeBERT, GraphCodeBERT, Hybrid bắt cặp từng từ trên vector encoder.
\item ASTNN bắt cặp câu lệnh trên vector RvNN.
\item CPG-GNN bắt cặp dòng trên vector node GAT.
\item \colorbox{evyellow!30}{Vàng} lẻ là token model dựa vào để cho điểm
  nhưng không ghép được cặp, không màu là phần còn lại.
\end{itemize}

\subsection{Cặp clone: 355523 với 355575}
Hai file giống hệt nhau.
Cả 5 phương pháp kết luận đúng: CodeBERT 1,000, GraphCodeBERT 1,000,
ASTNN 0,985, CPG-GNN 0,962, Hybrid 1,000, đều trên ngưỡng nên là clone.
\input{../workspace/demo_clone5.tex}

\subsection{Cặp không clone: bài 30 với 33}
Hai bài khác nhau của POJ-104.
Cả 5 phương pháp kết luận đúng: CodeBERT 0,000, GraphCodeBERT 0,010,
ASTNN 0,178, CPG-GNN 0,294, Hybrid 0,000, đều dưới ngưỡng nên không clone.
\input{../workspace/demo_neg5.tex}

\section{Hạn chế}
\begin{itemize}[leftmargin=*,itemsep=0pt,topsep=2pt,parsep=0pt]
\item Hybrid đỉnh ngay epoch đầu rồi overfit, CPG-GNN thì val F1 còn tăng
  ở epoch cuối.
\item CodeBERT, GraphCodeBERT cắt đầu vào 256 token,
  CPG-GNN chặn 1500 node, hàm dài bị mất thông tin.
\end{itemize}

\section{Kế hoạch tiếp theo}
\begin{itemize}[leftmargin=*,itemsep=0pt,topsep=2pt,parsep=0pt]
\item Giảm LR khi fine-tune Hybrid, train thêm vài epoch cho CPG-GNN.
\item Thử đầu vào 512 token, nới giới hạn node của CPG-GNN.
\end{itemize}

\end{document}
